% Chapitre : le jumeau numérique (axe 0).
% Sources principales : v2/jumeau/ (code), v2/artifacts/a0_jumeau/results.json (run de jumeau.check du
% 2026-11-15, 02:03), v2/docs/lab/reports/2026-11-14-jumeau-axe0.md, les trois revues adversariales
% v2/docs/lab/reports/2026-11-14-revue-adversariale-axe0{,-v2,-v3}.md et v2/papers/axe0-jumeau/main.tex.
% Les valeurs anciennes qui circulent encore dans le dépôt sont relevées en annexe (ann:incoherences).
\begingroup\sloppy\setlength{\emergencystretch}{3em}% coupures des chemins longs, local à ce fichier

\section{Le jumeau numérique}\label{sec:jumeau-ref}\label{jum:sec}

Le jumeau numérique est le moteur de projection commun à tous les axes. Il porte un état démographique
explicite, des modules qui rendent des flux typés et des scénarios faits d'événements datés. Le code est dans
\texttt{v2/jumeau/} ; les critères d'acceptation sont dans \texttt{v2/jumeau/check.py}, et leur dernier run est
consigné dans \texttt{v2/artifacts/a0\_jumeau/results.json} (2026-11-15, 02:03). Ce fichier n'est pas versionné
(\texttt{.gitignore} exclut \texttt{v2/artifacts/}) : il se régénère par \texttt{jumeau.check}
(section~\ref{jum:criteres}). Sauf mention contraire, toutes les valeurs mesurées de ce chapitre viennent de ce
fichier, donc de la composition de référence A0. Le contrat a été décidé le 2026-11-14 et corrigé
après chaque revue (\texttt{v2/docs/lab/decisions.md}, entrées « Contrat du jumeau (axe 0) » et suivantes).

Le jumeau n'est pas un modèle nouveau. Sa composition de référence réécrit le stock-flux de la v1 et le
reproduit à l'écart nul (critère C1). Cet écart nul est \emph{construit} : il montre que le contrat ne déforme
rien, pas que le modèle est exact. La précision de la projection relève du banc de validation, pas de ce
chapitre.

Dans ce chapitre, C1 à C18 désignent les critères d'acceptation de \texttt{check.py}, et T son relevé de temps ; l'estimateur de
cohortes de l'axe~2, aussi appelé C1, n'y intervient pas.

% ---------------------------------------------------------------------------------------------------
\subsection{Formalisation de l'état}\label{jum:etat}

\paragraph{Ensembles d'indices.} On note :
\begin{itemize}
  \item $\mathcal{P}$ les séries de profession. Le run de référence en compte 7\mesure{} : sages-femmes,
        chirurgiens-dentistes, pharmaciens, masseurs-kinésithérapeutes, infirmiers libéraux, infirmiers
        salariés, et une série « spécialistes » de test héritée de la v1, sans validation (\texttt{PROFS},
        \texttt{v2/jumeau/check.py}) ;
  \item $\Sigma$ les secteurs : privé, public, et un secteur \texttt{bloques} créé vide par la régulation
        (section~\ref{jum:leviers}). La sortie « libérale » du projet est le secteur privé, tout employeur
        public exclu (décision du 2026-11-14, « Recadrage en recherche ») ;
  \item $\mathcal{T}_\ell$ les territoires d'un niveau $\ell \in \{\text{commune}, \text{dép.}, \text{rég.},
        \text{nat.}\}$, déclaré par l'état ;
  \item $\mathcal{S} = \{\text{femme}, \text{homme}\}$ et $\mathcal{A} = \{a_1, \dots, a_J\}$ les tranches
        d'âge, de largeurs $w_j$ non nécessairement égales (critère C10) ;
  \item $t \in \{t_0, \dots, 2050\}$ les années, où $t_0$ est l'origine de la projection.
\end{itemize}

\paragraph{État.} Une cellule est un quadruplet $c = (p, \sigma, \tau, s) \in \mathcal{C} \subseteq
\mathcal{P} \times \Sigma \times \mathcal{T}_\ell \times \mathcal{S}$. L'état à l'année $t$ est la matrice
\begin{equation}\label{jum:eq:etat}
  N(t) = \bigl(N_{c,a}(t)\bigr)_{c \in \mathcal{C},\, a \in \mathcal{A}} \in \mathbb{R}_{+}^{|\mathcal{C}| \times
  |\mathcal{A}|}.
\end{equation}
Le secteur est un \emph{axe} de l'état, pas un filtre appliqué en sortie : un entrant bloqué, ou une transition
privé $\leftrightarrow$ public, change de cellule au lieu de disparaître. À cet état s'ajoute un vivier de
formation national $V(t) \in \mathbb{R}_{+}^{|\mathcal{P}| \times K}$, indexé par profession et année d'étude
$k = 1, \dots, L_p$ (\texttt{v2/jumeau/modules/formation.py}).

\paragraph{Agrégation.} Le passage d'un niveau $\ell$ à un niveau $\ell'$ plus grossier est une matrice
$G^{\ell \to \ell'} \in \{0,1\}^{|\mathcal{T}_{\ell'}| \times |\mathcal{T}_\ell|}$ dont chaque colonne contient
exactement un 1. Une unité non rattachée, ou rattachée deux fois, est refusée (\texttt{v2/papers/axe0-jumeau/main.tex},
section « L'état »). La sortie étudiée par les autres chapitres est
\begin{equation}\label{jum:eq:sortie}
  Y_{p,\tau',t} = \sum_{\tau \,:\, G_{\tau'\tau} = 1} \;\sum_{s \in \mathcal{S}} \sum_{a \in \mathcal{A}}
  N_{(p,\,\text{privé},\,\tau,\,s),\,a}(t).
\end{equation}
Le critère C8 vérifie que l'agrégation commune $\to$ département $\to$ région $\to$ national est exacte (écart
interne $0$\mesure{} aux trois niveaux, sur les \num{13147}\mesure{} communes présentes dans CartoSanté 2025, sur \num{34917}).

\paragraph{Taille.} Au département, l'état des 7 professions compte \num{1410}\mesure{} cellules. À la taille
communale, l'état de 5 professions sur 2 secteurs compte \num{698212}\mesure{} cellules (critère C15). Ce
dernier état est \emph{ventilé} à partir du département par la population ; il n'existe pas encore d'état
communal observé et projetable (section~\ref{jum:limites}).

% ---------------------------------------------------------------------------------------------------
\subsection{Modules et flux}\label{jum:modules}

\paragraph{Déclaration.} Un module est une classe enregistrée par \texttt{@register} et découverte par
\texttt{discover()} (\texttt{v2/jumeau/module.py}). Il déclare un nom, un rang d'application, des paramètres et
les flux de l'année qu'il lit (\texttt{needs}) ou produit (\texttt{produces}). Chaque paramètre porte une valeur,
une unité, une source, un millésime, un statut (mesuré, sourcé, supposé ou non mesurable) et des bornes. Le
tableau~\ref{jum:tab:modules} donne les sept modules actuels.

\begin{table}[h]
\centering\small
\begin{tabular}{lrlp{7.2cm}}
\toprule
Module & Rang & \texttt{needs} & Rôle et type de flux \\
\midrule
\texttt{formation}      & 5  & --- & Vivier national : passage $k \to k+1$ avec abandon, diplômés, places (Pool) \\
\texttt{vieillissement} & 10 & --- & Passage d'une tranche d'âge à la suivante (Shift) \\
\texttt{sorties}        & 20 & --- & Rétention nette ajustée ; leviers retraite, cumul, départs précoces (Rate) \\
\texttt{entrees}        & 30 & --- & Entrées nettes ajustées ; levier quotas (Add) \\
\texttt{lien\_formation} & 35 & \texttt{entrees}, \texttt{formation} & Part formée en France des entrées
                          $\times\, k(t)$ (Add) \\
\texttt{etrangers}      & 36 & \texttt{entrees} & Part des diplômés à l'étranger des entrées (Add) \\
\texttt{regulation}     & 40 & \texttt{entrees}, \texttt{sorties} & Plafond « entrées $\le$ sorties » en zone
                          régulée ; produit \texttt{regulation}, \texttt{regulation\_sortie} (Move, Add) \\
\bottomrule
\end{tabular}
\caption{Modules du jumeau, dans l'ordre d'application (\texttt{v2/jumeau/modules/*.py}).}
\label{jum:tab:modules}
\end{table}

\paragraph{Flux typés.} Un module ne modifie jamais l'état. Il rend des flux, appliqués par le moteur :
\begin{itemize}
  \item \textbf{Shift} (vieillissement) : $N_{c,a_j} \leftarrow (1 - f_j)\,N_{c,a_j} + f_{j-1}\,N_{c,a_{j-1}}$ ; la
        tranche ouverte garde tout ($f_J = 0$). Avec des tranches quinquennales, $f_j = 1/5$\suppose{} (âges
        supposés uniformes dans la tranche) ; avec des largeurs variables, $f_j = 1/w_j$ (C10 : fractions
        $(0{,}067 ; 0{,}1 ; 0{,}2 ; 0{,}2 ; 0)$\mesure{}, total conservé) ;
  \item \textbf{Rate} : $N \leftarrow s \odot N$, où $(1 - s) \odot N$ va dans un compte de sortie du registre ;
  \item \textbf{Add} : $N \leftarrow N + I$ ;
  \item \textbf{Move} : transfert d'une masse d'une cellule vers une autre, à total constant ;
  \item \textbf{Pool} : mise à jour du vivier, avec des comptes de diplômés et d'abandons.
\end{itemize}

\paragraph{Moteur.} Soit $m_1, \dots, m_M$ les modules de la composition, triés par rang. Le moteur
\texttt{run()} (\texttt{v2/jumeau/engine.py}) applique les flux module par module, par séparation d'opérateurs :
\begin{equation}\label{jum:eq:moteur}
  N^{(0)} = N(t), \qquad N^{(i)} = \Phi_{m_i}\bigl(N^{(i-1)}, \mathrm{ctx}_t\bigr), \qquad
  N(t+1) = \max\bigl(0,\, N^{(M)}\bigr),
\end{equation}
où $\Phi_{m}$ applique les flux rendus par $m$ et $\mathrm{ctx}_t$ donne accès aux flux déclarés des modules
précédents de l'année. Le plancher à 0 est appliqué à la fin de l'année. Chaque flux est inscrit par cellule
dans un registre. Un chemin rapide (\texttt{ledger=False}) omet le registre ; il donne un état identique au
chemin complet (écart $0{,}0$\mesure{}, C15).

\paragraph{Stock-flux de référence.} La composition « tendanciel » ne contient que \texttt{vieillissement},
\texttt{sorties} et \texttt{entrees}. Avec $b = \mathrm{Shift}(N(t))$ et $f = 1/5$\suppose{} :
\begin{equation}\label{jum:eq:stockflux}
  N(t+1) = \max\bigl(0,\; s \odot b + k(t)\,(I + p\,h)\bigr), \qquad h = t - t_0,
\end{equation}
où $s \le 1$ est la rétention nette ajustée par cellule et âge, $I$ l'entrée nette annuelle, $p$ sa pente et
$k(t)$ le multiplicateur de filière (\texttt{v2/papers/axe0-jumeau/main.tex}, équation du stock-flux). $s$, $I$
et $p$ sont ajustés par la fonction \texttt{fit\_rates} de la v1 (\texttt{v2/sante/model.py}), qui sert
d'oracle figé ; les opérations flottantes sont celles de la v1, dans le même ordre.

\paragraph{Pont vers le banc.} \texttt{twin\_of(EstimatorV1)} enveloppe le jumeau derrière l'interface commune
\texttt{v2/sante/estimators.py}. \texttt{trajectory(until, scenario)} exécute les modules du scénario. Le
jumeau se compare donc aux autres estimateurs sur le même banc, sans traitement particulier.

% ---------------------------------------------------------------------------------------------------
\subsection{Contrat de dépendances}\label{jum:contrat}

Le contrat a deux volets (\texttt{v2/jumeau/engine.py}, critère C5) :
\begin{itemize}
  \item \textbf{À la construction}, \texttt{build()} vérifie que chaque flux lu (\texttt{needs}) est produit par
        un module placé avant, et qu'aucun flux n'a deux producteurs. Une composition « régulation sans
        \texttt{sorties} » lève une \texttt{ValueError} explicite\mesure{}, au lieu de bloquer en silence tous
        les entrants comme avant la première revue.
  \item \textbf{Pendant le run}, toute lecture et tout flux non déclarés sont refusés. Les flux Pool se lisent
        par \texttt{ctx.pool\_out} : un accès par \texttt{ctx.delta} lève une \texttt{TypeError}\mesure{}, et
        \texttt{ctx.pool\_out} rend les \num{4725}\mesure{} diplômés attendus.
\end{itemize}
Ajouter un module revient à déposer un fichier : le critère C3 dépose un module factice hors du paquet et
vérifie que \texttt{twin\_of} l'exécute (flux inscrit \num{4040}\mesure{} = attendu), que le tendanciel est
inchangé (écart $0{,}0$\mesure{}) et qu'un module inconnu est refusé. Une méthode qui n'est pas un stock-flux s'y
branche de la même façon, en rendant ses résultats sous forme de flux.

% ---------------------------------------------------------------------------------------------------
\subsection{Scénarios et événements datés}\label{jum:scenario}

\paragraph{Contenu.} Un scénario (\texttt{v2/jumeau/scenario.py}, \texttt{FORMAT = 1}) est un quadruplet
$\mathcal{Z} = (\Theta, \mathcal{E}, \mathcal{M}, \omega)$ :
\begin{itemize}
  \item $\Theta$ : des surcharges \texttt{module.param} ;
  \item $\mathcal{E}$ : des événements datés $e = (\pi, t_d, t_f, \mathrm{op}, \mathrm{sel}, v)$, où $\pi$ est un
        paramètre, $[t_d, t_f]$ la période, $\mathrm{op} \in \{\texttt{set}, \texttt{mul}, \texttt{add}\}$ et
        $\mathrm{sel}$ un sélecteur sur les axes de l'état (profession, territoire, secteur, sexe) ;
  \item $\mathcal{M}$ : une composition de modules ;
  \item $\omega$ : une graine, exposée aux modules par \texttt{ctx.rng}.
\end{itemize}
La valeur effective d'un paramètre $\pi$ pour une cellule $c$ et une année $t$ est obtenue en appliquant, dans
l'ordre, les événements actifs dont le sélecteur couvre $c$ :
\begin{equation}\label{jum:eq:evenement}
  \pi_{c}(t) = \mathrm{op}_{e_n} \circ \cdots \circ \mathrm{op}_{e_1}\bigl(\pi^{\Theta}\bigr), \qquad
  \{e_1, \dots, e_n\} = \{e \in \mathcal{E} : \pi_e = \pi,\ t_d \le t \le t_f,\ c \in \mathrm{sel}_e\}.
\end{equation}

\paragraph{Composition avec la base.} Un scénario est \emph{posé sur} le scénario de base de l'estimateur par
\texttt{over()} : ses paramètres priment, ses événements s'ajoutent. Le critère C11 vérifie que
\texttt{trajectory(Scenario())} rend la même valeur que \texttt{predict} (\num{102987.2}\mesure{}).

\paragraph{Validation.} Avant le run sont refusés : les noms de paramètres, d'opérations et d'axes inconnus ; les
valeurs de sélecteur absentes de l'état (par exemple \texttt{'dentiste'} au lieu de \texttt{'dentistes'}) ; les
valeurs hors des bornes déclarées, y compris sur la valeur \emph{composée} par plusieurs événements et sur les
paramètres ajustés (C6, C13). Le scénario fait un aller-retour JSON sans perte (écart $0{,}0$\mesure{}), y compris
sa composition de modules après \texttt{over()} (C13).

\paragraph{Compositions disponibles.} « Tendanciel » : les trois modules du stock-flux v1. « REG » : le
tendanciel et \texttt{regulation}. « \texttt{LEVIERS} » : les sept modules (\texttt{scenario.py}) ; aux valeurs
par défaut, elle est identique au tendanciel (écart $0{,}0$\mesure{}, C14).

\paragraph{Délais.} Un multiplicateur de quotas décidé en $t_0$ agit à partir de $t_0 + L_p + 1$, où $L_p$ est la
durée des études (convention v1). Les durées sont reprises de la v1 (\texttt{v2/sante/levers\_v1.py},
\texttt{LAG}) : 10 ans pour les spécialistes, 6 pour les dentistes et les pharmaciens, 5 pour les sages-femmes et
les kinésithérapeutes, 3 pour les infirmiers ; elles sont supposées\suppose{}.

% ---------------------------------------------------------------------------------------------------
\subsection{Les leviers}\label{jum:leviers}

Le jumeau expose 12\mesure{} leviers en quatre familles : former, attirer, installer, retenir (C14). Chacun est
un paramètre d'un module, réglable par surcharge ou par événement daté. Le critère C14 vérifie l'identité au
tendanciel, l'effet non nul de chaque levier, son délai et sa portée régionale ; 9 mutations qui retirent l'effet
d'un levier le font échouer (\texttt{v2/docs/lab/reports/2026-11-14-jumeau-axe0.md}, version 5). Les effets
ci-dessous sont mesurés en 2050 par \texttt{jumeau.check}, sur l'estimateur \texttt{StockFlowLocal} des
dentistes et des pharmaciens (origine 2026), sauf mention contraire. Ce sont des tests du mécanisme, pas des
résultats de politique publique.

\paragraph{Former.}
\begin{itemize}
  \item \textbf{Quotas} (\texttt{entrees.multiplicateur}, bornes $[0 ; 3]$) : les entrées nettes sont multipliées
        par $m$ après le délai $L_p$. Le levier accepte un sélecteur régional. Quotas $\times 1{,}1$ : écart
        $4{,}5\cdot10^{-13}$\mesure{} avec le chemin v1 (\texttt{v2/sante/levers\_v1.project\_flow}, C7).
        Quotas $\times 1{,}2$ en Auvergne-Rhône-Alpes (84) : $+11{,}9$\mesure{}\,\% dans la région, 0\mesure{}
        ailleurs.
  \item \textbf{Capacité de formation} (\texttt{formation.capacite}) : les places de première année sont
        multipliées par $c$. Le vivier avance d'une année d'étude par an avec un abandon $\alpha_p$ ; la dernière
        année donne les diplômés $g_p(t)$. Un vivier de référence avance en parallèle aux valeurs par défaut et
        donne $g^{\mathrm{ref}}_p(t)$. Le module \texttt{lien\_formation} ajoute alors aux entrées positives
        $I^{+}$ la quantité
        \begin{equation}\label{jum:eq:lien}
          \bigl(1 - \varphi_p\bigr)\, I^{+}\, \bigl(k_p(t) - 1\bigr), \qquad k_p(t) = g_p(t) / g^{\mathrm{ref}}_p(t),
        \end{equation}
        où $\varphi_p$ est la part des diplômés à l'étranger (ci-dessous). $k_p = 1$ exactement au tendanciel.
        Capacité dentaire $\times 1{,}2$ : effet nul\mesure{} jusqu'en 2032, $+0{,}51$\mesure{}\,\% en 2033,
        $+6{,}3$\mesure{}\,\% en 2050 ; aucun effet sur les pharmaciens (0\mesure{}). DD76 donne $\pm 9$\,\% pour
        $\pm 20$\,\% de quotas \citep{drees2021dd76} ; l'écart vient de la part étrangère (voir « Limites »).
  \item \textbf{Abandon en formation} (\texttt{formation.abandon}) : taux annuel $\alpha_p$. Pour les infirmiers,
        $\alpha = 1 - 0{,}718^{1/3} = 0{,}1045$\derive{} par an, dérivé d'un rendement diplômés/inscrits mesuré
        (DREES, jeu 491). Kinésithérapeutes $0{,}016$\derive{}, sages-femmes $0{,}033$\derive{} ; autres
        professions $0$\suppose{}. Perte divisée par deux : diplômés $\times 1{,}196$\mesure{}, infirmiers libéraux
        $+14{,}1$\mesure{}\,\% en 2050.
\end{itemize}

\paragraph{Attirer.}
\begin{itemize}
  \item \textbf{Diplômés à l'étranger} (\texttt{etrangers.multiplicateur}, bornes $[0 ; 3]$) : le RPPS ne
        distingue pas le lieu de diplôme, donc les entrées ajustées contiennent déjà ces arrivées. Le module
        isole leur part au lieu d'ajouter un flux (\texttt{v2/jumeau/modules/etrangers.py}) :
        \begin{equation}\label{jum:eq:etrangers}
          \varphi_p = \min\Bigl(1,\; F_p \big/ \textstyle\sum I^{+}_p\Bigr), \qquad
          \text{ajout} = \varphi_p\, I^{+}\, (m - 1),
        \end{equation}
        où $F_p$ est le flux annuel de diplômés à l'étranger : dentistes \num{1300}\derive{}, kinésithérapeutes
        \num{2200}\derive{}, pharmaciens 240\derive{}, sages-femmes 80\derive{}, infirmiers 400, sourcé
        avec une confiance faible (DREES, ER 1319), dont seule la borne haute de \num{1100}\suppose{} est supposée
        (\texttt{FLUX} ; étiquettes de la décision du 2026-11-15). L'ajout entre dans la tranche d'âge qui contient
        28~ans, âge d'entrée des diplômés à l'étranger donné par DD76 (tableaux 12 et 13) ; pour les
        kinésithérapeutes, l'ER 1075 donne 26~ans, dans la même tranche (C17). Étrangers $\times 2$ : part dentaire $\varphi = 0{,}47$\mesure{}, dentistes
        $+37{,}2$\mesure{}\,\%, pharmaciens $+6{,}1$\mesure{}\,\% en 2050 ; un flux nul ne produit aucun effet
        (0\mesure{}).
\end{itemize}

\paragraph{Installer.} Le module \texttt{regulation} représente la règle conventionnelle « une installation
pour un départ » dans les zones que le zonage conventionnel classe en densité la plus élevée (\texttt{v2/jumeau/modules/regulation.py}). Pour chaque cellule
privée d'une zone régulée et chaque année :
\begin{equation}\label{jum:eq:regulation}
  x = \max\bigl(E - D,\, 0\bigr)\,\bigl(1 - \epsilon_p\bigr), \qquad B = \theta\, x,
\end{equation}
où $E$ et $D$ sont les entrées et les sorties de l'année, $\epsilon_p$ la part dérogatoire et $\theta$ la part
contraignante du plafond. Les bloqués $B$ se répartissent entre un territoire voisin non régulé (privé), le
salariat (secteur \texttt{bloques}) et la sortie. Cinq leviers agissent sur ce mécanisme :
\begin{itemize}
  \item \texttt{regulation.zones} : la liste des géocodes régulés, par profession ;
  \item \texttt{regulation.part\_zones} : la part des zones listées qui est retenue, dans l'ordre de la liste. Pour
        les kinésithérapeutes, $0 / 0{,}6 / 1$ déplacent $0 / 227 / 356$\mesure{} personnes la première année ;
  \item \texttt{regulation.theta} : $\theta = 0{,}80$\suppose{} [0,74 ; 0,87], estimé par la DiD N3 sur le zonage
        infirmier, dont le test de tendances parallèles échoue ($p = 0{,}043$) : valeur exploratoire, non
        validée ; C12 vérifie que $0{,}800$\mesure{} de l'excédent est bloqué ;
  \item \texttt{regulation.repartition} : par défaut 100\suppose{}\,\% vers un voisin, ce qui rend la régulation
        sans effet au département, par construction (C12 : écart $0{,}0$\mesure{}) ;
  \item \texttt{regulation.part\_exemptee} : la part $\epsilon_p$. Pour les dentistes spécialistes, la convention
        de 2023 (art.~35.5) prévoit une dérogation ; sa part parmi les entrants n'est pas mesurée et vaut
        0\suppose{} (décision du 2026-11-15).
\end{itemize}
Sur l'état réel à deux secteurs, avec des zones illustratives (départements 06, 13, 34, 66, 83 ;
\texttt{ZONES}, \texttt{check.py}), $\theta = 0{,}8$ et 100\,\% des bloqués vers le salariat, l'effectif privé des
zones varie de \num{-6711}\mesure{} kinésithérapeutes et de \num{-3484}\mesure{} dentistes en 2050. Le total ne change pas
(écart $0{,}0$\mesure{} sans plancher) : la régulation déplace, elle ne crée ni ne supprime (C2). Si les bloqués
partent vers le secteur public existant, qui a sa propre démographie, le total varie de \num{-375}\mesure{}
(kinésithérapeutes) et \num{+203}\mesure{} (dentistes).

\paragraph{Retenir.} Le module \texttt{sorties} agit sur les taux de sortie $e = 1 - s$
(\texttt{v2/jumeau/modules/sorties.py}).
\begin{itemize}
  \item \textbf{Report de la retraite} (\texttt{sorties.decalage\_retraite}, $\Delta$ ans, bornes $[0 ; 5]$),
        mécanisme dit B1 :
        \begin{equation}\label{jum:eq:retraite}
          e' =
          \begin{cases}
            e / (1 + \lambda e) & \text{tranches de 65~ans et plus,} \\
            e\, \max(0,\, 1 - \Delta/5) & \text{60--64~ans,} \\
            e\, \max(0,\, 1 - \Delta/10) & \text{55--59~ans,}
          \end{cases}
        \end{equation}
        les tranches à solde net positif ($e = 0$) restant inchangées. Avec une seule tranche ouverte de
        65~ans et plus, $\lambda = \Delta$ et le séjour moyen passe de $1/e$ à $1/e + \Delta$. Avec plusieurs
        tranches ($65$--$69$, $70$--$74$, $75+$), $\lambda$ est calculé par cellule pour que le temps espéré en
        65+ d'une personne qui entre en 65--69 augmente de $\Delta$ exactement ; ce temps est l'espérance d'une
        chaîne de Markov dont les transitions sont $f$ et $s$, dans l'ordre du moteur. La montée en charge est
        de 3\suppose{}~mois par an à partir de l'année qui suit la décision. Mesures (C7, C13) : séjour en 65+
        allongé de $+1{,}000$ / $+2{,}000$ / $+5{,}000$\mesure{}~an pour $\Delta = 1 / 2 / 5$ (écart maximal
        $3{,}6\cdot10^{-14}$\mesure{}) ; avec trois tranches de 65~ans et plus, écart de séjour
        $9{,}8\cdot10^{-15}$\mesure{}. Effet en 2050 : dentistes $+1{,}60$\mesure{}\,\% ($\Delta = 1$) et
        $+3{,}22$\mesure{}\,\% ($\Delta = 2$) ; kinésithérapeutes $+1{,}55$\mesure{}\,\% et infirmiers libéraux
        $+3{,}02$\mesure{}\,\% pour $\Delta = 1$.
  \item \textbf{Cumul emploi-retraite} (\texttt{sorties.cumul\_multiplicateur}, $k$) :
        $e'_{65+} = e_{65+}\, \max\bigl(0,\, 1 - \kappa_p (k - 1)\bigr)$, où $\kappa_p$ est la part des libéraux en
        cumul : infirmiers et kinésithérapeutes 0,015, sages-femmes 0,008, dentistes 0,047, pharmaciens
        0,018 au 30/06/2015, sourcés (COR 2016, doc.~11, annexe~2). Appliquer la part de la CARPIMKO aux
        kinésithérapeutes et aux infirmiers est supposé\suppose{}.
  \item \textbf{Départs précoces} (\texttt{sorties.multiplicateur\_precoces}, $m$) : $e'_{a < 60} = m\, e_{a<60}$.
        Aucune source ne chiffre l'effet d'une politique sur ce taux ; le mécanisme est supposé\suppose{}.
        $m = 0{,}5$ : $+5{,}6$\mesure{}\,\% en 2050 (dentistes et pharmaciens réunis).
\end{itemize}

% ---------------------------------------------------------------------------------------------------
\subsection{Critères d'acceptation}\label{jum:criteres}

Les critères sont écrits pour pouvoir échouer. Pour chacun, \texttt{check.py} garde la valeur mesurée par une
revue adversariale sur le code d'avant (champ \texttt{avant} de \texttt{results.json}) ; c'est la preuve que le
test sait échouer. Commande, depuis la racine du dépôt :
\begin{verbatim}
PYTHONPATH=v2 uv run python -m jumeau.check
\end{verbatim}
Le code de sortie vaut 1 si un critère échoue. Le run du 2026-11-15 (02:03) passe les 18\mesure{} critères sur
18 (tableau~\ref{jum:tab:criteres}). La ligne T est un relevé de temps, pas un test : son verdict est codé
\texttt{True} (\texttt{check.py}, l.~808--810) et elle ne peut pas échouer. Son extrapolation à 5 professions,
$\approx 2{,}1$~s\derive, dépasse d'ailleurs la cible de 2~s ; seul C15 vérifie la cible, sur le chemin rapide.
Plusieurs critères lisent à l'exécution des fichiers bruts non commités : C16 lit le jeu DREES 491
(\texttt{data/raw/drees/}, ouvert, à télécharger). C5, C9 et la variante filière (C1, C11) lisent les inscrits
en odontologie et en pharmacie de SISE (données ouvertes, \texttt{v2/data/processed/sise\_sante.parquet}, commité),
qui remplacent depuis le 15 novembre 2026 un fichier fourni sous accord de confidentialité ; SISE le reproduit à
l'identique pour les années d'études 2 à 5\mesure, mais s'arrête à la rentrée 2024 (chapitre~\ref{sec:repro-ref}).
Les valeurs de C5, C9 et C11 du tableau sont celles du run avec SISE : vivier de 5\,064 diplômés\mesure (C5),
4\,473\mesure (C9), dentistes 2050 en variante filière 85\,252\mesure au lieu de 102\,987 avec l'ancien fichier
(C11, $-17$\,\%\derive, effet de la rentrée 2025 manquante, que la variante prolonge jusqu'en 2050).

{\small
\begin{longtable}{@{}lp{5.2cm}p{8.0cm}@{}}
\caption{Critères d'acceptation du jumeau : dernier run de \texttt{jumeau.check}
(\texttt{v2/artifacts/a0\_jumeau/results.json}, 2026-11-15, 02:03). Toutes les valeurs sont
mesurées\mesure{}, sauf mention.}\label{jum:tab:criteres}\\
\toprule
\# & Ce qui est testé & Valeur mesurée \\
\midrule
\endfirsthead
\toprule
\# & Ce qui est testé & Valeur mesurée \\
\midrule
\endhead
\bottomrule
\endfoot
C1 & Reproduction du stock-flux v1 (5 variantes, banc K1 (backtest départemental de la v1, ancien périmètre libéral + mixte) et origine récente $\to$ 2050) &
  écart max $0{,}0$ sur 140 comparaisons de \texttt{predict} ; $0{,}0$ sur \num{355248} valeurs
  cellule $\times$ âge $\times$ année, 7 professions \\
C2 & Régulation exécutée par \texttt{twin\_of}, état réel à 2 secteurs, sans créer ni supprimer &
  privé des zones \num{-6711} (MK), \num{-3484} (CD) ; $\Delta$ total sans plancher $0{,}0$ ; vers le
  public \num{-375} / \num{+203} \\
C3 & Module factice déposé hors du paquet & flux \num{4040} = attendu ; tendanciel inchangé ($0{,}0$) ; module
  inconnu refusé \\
C4 & Régulation vers un secteur vide, mêmes taux : déplacement pur & $\Delta$ total 2050 $0{,}0$ (avec plancher
  $0{,}000$) ; 202 cellules en repli \\
C5 & Dépendances déclarées et vérifiées & dépendance manquante, double producteur, lecture non déclarée :
  refusés ; Pool par \texttt{ctx.delta} : refusé ; par \texttt{ctx.pool\_out} : \num{4725} diplômés \\
C6 & Scénario validé et sérialisable & faute de frappe, opération, axe, sélecteur \texttt{'dentiste'} ou
  \texttt{6}, quotas $\times 110$ : refusés ; aller-retour JSON $0{,}0$ \\
C7 & Leviers quotas et retraite contre leur référence & quotas $\times 1{,}1$ : $4{,}5\cdot10^{-13}$ avec la v1 ;
  séjour 65+ $+1{,}000 / +2{,}000 / +5{,}000$ an ; montée $[0{,}25 ; 0{,}5 ; 0{,}75 ; 1]$ \\
C8 & Agrégation commune $\to$ national ; écart au RPPS & écart interne 0 (dép., rég., nat.) ; commune inconnue
  refusée ; CartoSanté contre RPPS 2025 : CD $-2{,}8$\,\%, MK $-5{,}4$\,\% \\
C9 & Aucun état partagé entre deux runs & diplômés 2028 : vivier $\times 1$ \num{4253}, $\times 2$ \num{8505} =
  instances neuves \num{8505} \\
C10 & Vieillissement à largeurs variables & fractions $(0{,}067 ; 0{,}1 ; 0{,}2 ; 0{,}2 ; 0)$ ; total
  $430 \to 430$ \\
C11 & Scénario posé sur la base (\texttt{over}) & \texttt{predict} = \texttt{trajectory(Scenario())} =
  \num{102987.2} ; quotas $\times 1{,}1$ : $+6{,}53$\,\% (CD, $L_p = 6$\suppose{}) \\
C12 & Paramètres sourcés de l'axe~1 câblés & $\theta$ 0,8 $\to$ 0,800 des bloqués ; dérogation 50\,\% $\to$
  0,500 ; 100\,\% voisine au département $\to$ écart 0 ; 100\,\% sortie, année 1 : $\Delta$ \num{-355.8} =
  compte ; abandon IDE 0,1045\derive{} ; cumul 0,015 / 0,008 / 0,047 \\
C13 & Défauts de la troisième revue & bornes composées et ajustées, NaN, rétention $\times 1{,}3$ : refusés ;
  JSON + \texttt{over} garde la régulation ; règle « entrées $\le$ sorties » : écart $7{,}1\cdot10^{-15}$ à une
  formule indépendante ; séjour à 3 tranches $\ge 65$ : $9{,}8\cdot10^{-15}$ \\
C14 & 12 leviers branchés (identité, effet, délai, région) & identité $0{,}0$ ; capacité CD $\times 1{,}2$ : 0 en
  2032, $+0{,}51$\,\% en 2033, $+6{,}3$\,\% en 2050 ; IFSI : $+14{,}1$\,\% ; étrangers $\times 2$ : $+37{,}2$\,\%
  CD ; précoces $\times 0{,}5$ : $+5{,}6$\,\% ; quotas région 84 : 0 hors région \\
C15 & Vitesse : 5 professions à la commune $< 2$~s, cible supposée du profil\suppose{} & \SI{1.33}{\second} réel, \SI{1.28}{\second}
  CPU (médianes de 3 runs ; \num{698212} cellules ; 24 ans ; charge 4,5) ; chemin rapide = complet \\
C16 & Vivier réel (DREES 491) branché & identité en mode levier $0{,}0$ ; diplômés IDE \num{25385} contre
  \num{25249} (2024), \num{26005} contre \num{26020} (2025) ; MAPE vivier / persistance 2019--2025 : IDE 8,4 /
  3,2\,\%, MK 5,9 / 3,1\,\%, SF 4,4 / 4,1\,\% \\
C17 & Diplômés à l'étranger entrés à 28~ans & ajout dans la seule tranche qui contient 28~ans (libellée
  25+) ; hors tranche $0{,}0$ ; ajout année 1 = \num{1300} \\
C18 & Infirmiers par secteur (axe~4) & \num{1100} dép. $\times$ années, écart max $7{,}3\cdot10^{-12}$ à SNDS +
  BTS ; 2023 : privé \num{277788}, public \num{329856} ; privé 2050 (tendanciel) \num{336422} \\
T & Relevé de temps jusqu'en 2050 (verdict codé \texttt{True}, ne peut pas échouer) & \SI{16}{\milli\second}
  (\texttt{run()} direct d'A0, dép., 7 séries) ; \SI{0.42}{\second} (taille commune, 1 profession $\times$ 2
  secteurs) ; $\approx$ \SI{2.1}{\second}\derive{} extrapolé à 5 professions, au-dessus de la cible de 2~s \\
\end{longtable}
}

Trois remarques sur la portée de ces critères.
\begin{itemize}
  \item C1, C3, C4, C8 (partie interne) et C10 vérifient des propriétés du contrat, en partie vraies par
        construction ; la première revue l'a relevé (section~\ref{jum:revues}). Elles ne disent rien de
        l'exactitude des projections.
  \item C8 publie un écart externe, non corrigé : CartoSanté 2025 compte \num{84946}\mesure{} kinésithérapeutes
        contre \num{89766}\mesure{} au RPPS ($-5{,}4$\mesure{}\,\%, dont \num{-2084}\mesure{} dans les
        Bouches-du-Rhône, où Marseille est absente) ; la MAPE départementale vaut 4,2\mesure{}\,\% (CD) et
        4,7\mesure{}\,\% (MK).
  \item C16 contient un résultat négatif : sur les diplômés 2019--2025, la persistance prévoit mieux que le
        vivier pour les trois professions testées. Sur le stock d'infirmiers libéraux 2018 $\to$ 2023, le mode
        vivier fait 7,43\mesure{}\,\% de MAPE départementale contre 7,82\mesure{}\,\% pour le stock-flux local,
        sur une seule origine : c'est exploratoire.
\end{itemize}

% ---------------------------------------------------------------------------------------------------
\subsection{Les trois revues adversariales}\label{jum:revues}

Chaque revue a été conduite par des agents qui n'ont modifié aucun fichier, sur un état daté du code. Un défaut
n'est retenu que s'il résiste à une tentative de réfutation. Chaque défaut corrigé est couvert par un critère, et
une mutation qui le réintroduit fait échouer ce critère (\texttt{v2/docs/lab/reports/2026-11-14-jumeau-axe0.md},
section « Mutations »).

\paragraph{Revue 1} (\texttt{v2/docs/lab/reports/2026-11-14-revue-adversariale-axe0.md}). Neuf agents, trois angles
(les tests peuvent-ils échouer ; justesse de la conservation, de la régulation et de l'agrégation ; contrat). Le
jumeau annonçait alors 4 critères, dont 3 vrais par construction. Six défauts confirmés sur six soumis, dont
quatre bloquants et deux importants, plus trois constats non vérifiés :
\begin{itemize}
  \item \texttt{trajectory} ignorait \texttt{scenario.modules} ; le test plug-in ne pouvait pas échouer
        $\Rightarrow$ C3 passe par \texttt{twin\_of}, module inconnu refusé ;
  \item un secteur public sans sorties ($s = 1$) créait $+822$\mesure{} kinésithérapeutes en 2050 $\Rightarrow$
        règle de repli sur la cellule privée (supposée), C4 ;
  \item dépendances implicites $\Rightarrow$ \texttt{needs} / \texttt{produces} vérifiés par \texttt{build()}, C5 ;
  \item scénario non validé et non sérialisable $\Rightarrow$ validation des noms, JSON versionné, C6 ;
  \item le levier retraite était comparé à une copie de lui-même ; les quotas agissaient un an trop tôt
        ($+226$\mesure{} dentistes) $\Rightarrow$ C7 contre le vrai chemin v1 ;
  \item l'agrégation était vérifiée de façon circulaire $\Rightarrow$ C8 publie l'écart au RPPS.
\end{itemize}
Les constats non vérifiés (vieillissement codé à $1/5$, vivier \texttt{\_intake} persistant d'un run à l'autre,
conservation qui est une identité) ont mené à C9 et C10. Le calibrage de \texttt{fit\_rates} à $1/5$ reste
ouvert.

\paragraph{Revue 2} (\texttt{v2/docs/lab/reports/2026-11-14-revue-adversariale-axe0-v2.md}). Six constats sur six
confirmés, soit quatre défauts bloquants distincts :
\begin{itemize}
  \item le masque retraite touchait la tranche ouverte 65+ : séjour infini pour $\Delta \ge 5$, effet
        surestimé de $+41$\mesure{}\,\% à $\Delta = 2$ (dentistes : \num{3137} contre \num{2229} après
        correctif). Le code contredisait une décision inscrite $\Rightarrow$ mécanisme B1
        (équation~\ref{jum:eq:retraite}), C7 contre le séjour attendu ;
  \item la régulation faisait encore varier le total ($+254$\mesure{} dentistes, $-474$\mesure{} kinés)
        $\Rightarrow$ destination \texttt{bloques} qui conserve le total, C2 ;
  \item \texttt{trajectory(scenario)}, introduit par la correction précédente, écrasait le scénario de base :
        quotas $\times 1{,}1 \to -15{,}2$\mesure{}\,\% $\Rightarrow$ \texttt{over()}, C11 ;
  \item valeurs de sélecteur et plages non validées : quotas $\times 110 \to +$\num{5085033}\mesure{} dentistes
        $\Rightarrow$ bornes, C6.
\end{itemize}
Conséquence enregistrée : les chiffres retraite et régulation du prototype ont été déclarés non citables et les
leviers suspendus jusqu'à la revue suivante (\texttt{decisions.md}, 2026-11-14, « Le jumeau ne passe pas la
seconde revue »).

\paragraph{Revue 3} (\texttt{v2/docs/lab/reports/2026-11-14-revue-adversariale-axe0-v3.md}). Aucun défaut bloquant
dans le jumeau. Le seul défaut bloquant portait sur le prototype (\texttt{IndexError} à la réactivation de la
régulation, 606 cellules contre 404) et a été corrigé par un masque par trajectoire. Défauts importants corrigés et
couverts par C13 :
\begin{itemize}
  \item bornes contournables par empilement d'événements ($\times 243 \to +14\,317$\mesure{}\,\%) et par NaN ;
  \item rétention ajustée supérieure à 1 possible (rétention $\times 1{,}3 \to +$\num{127929}\mesure{} au lieu de
        \num{-2588}) ;
  \item l'aller-retour JSON avec \texttt{over} perdait la composition ;
  \item séjour différent de $+\Delta$ avec plusieurs tranches de 65~ans et plus $\Rightarrow$ $\lambda$ par cellule.
\end{itemize}
Les trois constats non vérifiés (règle « entrées $\le$ sorties » non testée, entrées négatives déplacées vers les
bloqués, événement sur les zones qui plantait) sont couverts depuis par C13. La revue mesure une conservation à
0,5\mesure{}~personne près en 2050 (CD, MK, SF). Les leviers retraite et régulation ont été réactivés, avec des
bornes côté serveur (\texttt{decisions.md}, 2026-11-14, « Troisième revue »).

\paragraph{Ce que les revues ne couvrent pas.} Les agents ont relu un état daté du code ; les constats non
réfutés n'ont pas tous été vérifiés par une mesure, et aucune revue n'a porté sur le code après C14--C18
(leviers branchés, vitesse, vivier réel, diplômés à l'étranger, infirmiers par secteur). Une quatrième revue coûte
environ 3\suppose{}~min de machine (\texttt{v2/docs/lab/reports/2026-11-14-revue-adversariale-axe0-v2.md}, questions
ouvertes).

% ---------------------------------------------------------------------------------------------------
\subsection{Temps de recalcul}\label{jum:temps}

Le profil du projet fixe une cible de recalcul inférieure à 2~s\suppose{} (\texttt{docs/lab/profile.md} ; seuil
de C15). Mesures du run de référence (\texttt{results.json}, bloc \texttt{temps} et critère C15) :
\begin{itemize}
  \item \textbf{Département}, 7 séries de profession (\num{1410} cellules), jusqu'en 2050 :
        \SI{16}{\milli\second}\mesure{} par scénario, médiane de plusieurs répétitions de \texttt{run()} direct
        d'A0 (sans le cache de \texttt{trajectory}). Ce n'est pas le temps servi : un scénario à tous leviers
        servi avec C1 prend 280 à 330\,ms\mesure{} (\texttt{decisions.md}, « C1 servi par la carte et
        l'assistant ») ;
  \item \textbf{Taille communale répliquée}, stock-flux en 3 modules : \SI{0.22}{\second}\mesure{} pour 1
        profession $\times$ 1 secteur (\num{70094} cellules) ; \SI{0.42}{\second}\mesure{} pour 1 profession
        $\times$ 2 secteurs (\num{140188} cellules) ; $\approx$ \SI{2.1}{\second}\derive{} extrapolé
        linéairement à 5 professions ;
  \item \textbf{C15}, 5 professions $\times$ 2 secteurs à la commune (\num{698212} cellules, 24 ans), chemin
        rapide : \SI{1.33}{\second}\mesure{} réel et \SI{1.28}{\second}\mesure{} CPU, médianes de 3 runs, charge
        machine 4,5\mesure{}. Avant optimisation : \SI{3.99}{\second}\mesure{}. D'après le profil fait avant
        optimisation, le temps se répartissait entre le registre (1,1\mesure{}~s), les passes (1,0\mesure{}~s),
        les sorties (0,8\mesure{}~s) et les entrées (0,5\mesure{}~s) (\texttt{v2/papers/axe0-jumeau/main.tex}).
\end{itemize}
Le temps à la commune dépend de la charge de la machine : des mesures antérieures sous charge ont dépassé
2~s (annexe~\ref{ann:incoherences}). La marge sous la cible est donc d'environ un facteur 1,5\derive{} sur une
machine peu chargée, et n'est pas établie sous charge. L'extrapolation T ($\approx 2{,}1$~s) et la mesure directe
C15 ($1{,}33$~s) ne portent pas sur la même composition : T prolonge un état répliqué à un secteur par
profession, C15 mesure le chemin rapide. Le prototype de carte, au département, mesure de son côté une réponse
perçue p50 de \SI{53}{\milli\second}\mesure{} (\texttt{v2/docs/lab/reports/2026-11-14-prototype-carte.md}).

% ---------------------------------------------------------------------------------------------------
\subsection{Limites}\label{jum:limites}

\begin{itemize}
  \item \textbf{C1 ne valide pas le modèle.} L'écart nul est construit. La précision reste celle du stock-flux v1
        au banc K1 : 8,08\mesure{}\,\% de MAPE départementale, sur l'ancien périmètre libéral + mixte
        (\texttt{v2/docs/lab/reports/2026-11-14-jumeau-axe0.md}, menaces).
  \item \textbf{Calibrage à $1/5$.} \texttt{fit\_rates} suppose des tranches quinquennales ; il n'est pas adapté
        aux tranches communales ni aux largeurs variables que le moteur accepte (C10). Ouvert depuis la première
        revue.
  \item \textbf{Plancher à 0.} Il crée 0,02 à 0,43\mesure{}\,\% de $N(2050)$ selon la profession ; il n'est pas
        inscrit comme un flux.
  \item \textbf{Ordre des modules.} La séparation d'opérateurs dépend de l'ordre d'application ; la sensibilité
        à cet ordre n'est pas mesurée (environ 15\suppose{}~min).
  \item \textbf{Régulation.} La destination des entrants bloqués n'est chiffrée par aucune source ; la
        destination \texttt{bloques} conserve le total par construction, et la variante « vers le public »
        change le total de \num{-375} / \num{+203}\mesure{}. Le défaut (100\,\% voisine) rend la régulation sans
        effet au département ; à la commune, elle exige une table de voisinage, qui n'existe pas. Les zones des
        critères (06, 13, 34, 66, 83) sont illustratives, ce ne sont pas les zonages conventionnels.
  \item \textbf{Repli du public.} Une cellule publique sans historique emprunte la rétention de la cellule privée,
        ce qui est supposé\suppose{}.
  \item \textbf{Diplômés à l'étranger.} Les flux mélangent des définitions (primo-inscriptions, entrées au
        RPPS, autorisations) ; un âge unique d'entrée est appliqué à toutes les professions. La part étrangère
        dentaire ($\varphi = 0{,}47$\mesure{}) explique que la capacité $\times 1{,}2$ donne $+6{,}3$\,\% contre
        $\pm 9$\,\% dans DD76 : la concordance avec la DREES dépend d'un paramètre mal mesuré.
  \item \textbf{IFSI.} Le jumeau donne $+14{,}1$\mesure{}\,\% d'infirmiers libéraux en 2050 pour une perte
        divisée par deux ; l'équivalent en effectifs du chiffre DREES (ER 1319) est d'environ
        $+7{,}9$\derive{}\,\%. L'écart, d'environ un facteur 2, n'est pas expliqué.
  \item \textbf{Retraite.} Le repère DREES d'environ $+1$\,\% n'est pas comparable aux $+1{,}55$ à
        $+3{,}02$\mesure{}\,\% du jumeau : il porte sur tous les modes d'exercice, les plus de 60~ans, un plafond à
        67~ans et dès 2020.
  \item \textbf{Vivier.} La persistance prévoit mieux les diplômés 2019--2025 que le vivier (C16) ; le mode
        vivier n'est pas le défaut, et la capacité agit par un vivier relatif dont seul le rapport à la référence
        compte.
  \item \textbf{Échelle communale.} Les temps à la commune sont mesurés sur un état ventilé ou répliqué ;
        l'état communal observé n'existe pas, et l'écart CartoSanté / RPPS ($-5{,}4$\mesure{}\,\% pour les MK)
        n'est pas corrigé.
  \item \textbf{Revues.} Aucune revue adversariale n'a porté sur C14--C18. Les statuts « supposé » de
        $\theta$, de la répartition, de la dérogation et des durées $L_p$ restent ceux des fiches de l'axe~1.
\end{itemize}
\endgroup
